\section{Technical derivations}
\label{dist:app:technical}

\subsection{Lifetime yield from \texorpdfstring{$k$}{k} founders}
\label{dist:app:vk}

We derive \cref{dist:eq:vk}. Start from
$V_\infty^{(k)}=\int_0^\infty g_k(t)\,\dd t$ with $g_k$ from
\cref{dist:eq:gkflux}, and change variables from $t$ to $x=I(t)$: as $t$ runs
$0\to\infty$, $I$ runs $1\to b$, and since $\dd I/\dd t=-\delta J$ we have
$\dd t=-\dd x/(\delta J)$. The first part of $g_k$ gives
\[
\int_0^\infty\delta k K I^{k-1}\dd t
=\int_b^1 k\,\frac{K}{J}\,x^{k-1}\dd x
=\int_b^1 k\Bigl(1+\frac{2\lambda}{\delta}-\frac{2\lambda}{\delta}x\Bigr)
x^{k-1}\dd x,
\]
using $K/J=1+2\lambda(1-I)/\delta$. Integrating,
\[
=\Bigl(1+\frac{2\lambda}{\delta}\Bigr)(1-b^k)
-\frac{2\lambda k}{\delta}\,\frac{1-b^{k+1}}{k+1}
=(1-b^k)+\frac{2\lambda}{\delta}
\Bigl[\frac{1}{k+1}-b^k+\frac{k\,b^{k+1}}{k+1}\Bigr],
\]
where the second form collects terms over the common denominator $k+1$. The
second part of $g_k$ gives, with $J=-\frac{\lambda}{\delta}(I-a)(I-b)$,
\[
\int_0^\infty\delta k(k-1)J^2 I^{k-2}\dd t
=\int_b^1 k(k-1)J\,x^{k-2}\dd x
=k(k-1)\frac{\lambda}{\delta}\int_1^b (x-a)(x-b)x^{k-2}\dd x,
\]
on reversing the limits. Adding the two parts yields \cref{dist:eq:vk}. The
formula also holds at $k=1$, the integral term vanishing, and reduces to
$V_\infty=(1-b)[1+\lambda(1-b)/\delta]$ of \cref{dist:eq:vinf}. When $\mu=0$,
$b=0$ and the integral simplifies to
$k(k-1)\frac{\lambda}{\delta}\int_1^0(x-a)x^{k-1}\dd x
=k(k-1)\frac{\lambda}{\delta}\bigl[\frac{a-1}{k}-\frac{1}{k+1}\bigr]$, which
combines with the first part to give $k+\lambda/\delta$ --- the elementary
result of \cref{dist:eq:vkmuzero}, recovered analytically.

\subsection{Coefficient extraction for general initial conditions}
\label{dist:app:extract}

The general-$\mu$ route for the chained model of \cref{dist:sec:chained} runs
through the $k$-founder PGF. From \cref{dist:sec:pgf-k-founders}, in the
catastrophe convention,
\begin{equation}
G_k(z,t)=\Biggl[\frac{ab\gamma(t)+v(t)z}{1-\gamma(t)z}\Biggr]^k,
\qquad
\gamma(t)=\frac{1-w}{b-aw},
\qquad
v(t)=\frac{bw-a}{b-aw},
\label{dist:eq:gkgamma}
\end{equation}
with $w=\ee^{\theta t}$; note that $\gamma$ is the sliding ratio $P$ of
\cref{dist:eq:pdef}. Writing
$\Theta_m=\binom{k}{m}(ab)^{k-m}\gamma(t)^{k-m}v(t)^m$ and
$\Omega_n=\binom{k-1+n}{n}\gamma(t)^n$, the binomial expansion of the
numerator and the negative-binomial series of $(1-\gamma z)^{-k}$ give
$G_k(z,t)=\bigl(\sum_{m=0}^k\Theta_m z^m\bigr)\bigl(\sum_{n\ge0}\Omega_n
z^n\bigr)$, so that the state-probability coefficients are, with $H$ here
denoting the power of $z$ rather than the rupture state,
\begin{equation}
[z^H]G_k=
\sum_{i=0}^{\min(H,k)}
\Theta_{k-i}\,\Omega_{i+\max(0,H-k)}.
\label{dist:eq:zHGk}
\end{equation}
The chained burst-size laws then require integrals of the form
$\int_0^\infty\gamma(t)^n v(t)^m\dd t$ for arbitrary $n,m\in\mathbb{N}$.
Binomial-expanding the two numerators,
\begin{equation}
(1-w)^n(bw-a)^m=\sum_{i=0}^{n+m}C_i\,w^{i}
\label{dist:eq:cexpand}
\end{equation}
for explicit constants $C_i$, and the problem reduces to the elementary
building blocks
\begin{equation}
\int\frac{w^{i}}{(b-aw)^j}\dd t
=-\frac{i\,w^{i}\,(b-aw)^{1-j}}{b\,\theta}\,
\Fhyp\Bigl(1,(1+i)-j;1+i;\frac{aw}{b}\Bigr),
\label{dist:eq:hypantideriv}
\end{equation}
whose integral identity is established in
\mextref{m:app:hyp}{the hypergeometric identity of \ChM}, at
\mextref{m:prop:hypIden}{that proposition} and
\mextref{m:eq:hypIden}{its displayed identity}. The burst-size law for the
$k$th rupture then follows by summing over the initial-condition
distribution. For $\mu=0$ this entire machinery collapses to the geometric
decomposition of \cref{dist:sec:key-observation}; for $\mu>0$ it remains the
available route, and the joint laws are the open problem of
\cref{dist:sec:open-problems}. One step here is not routine: the interchange
of summation and integration in \cref{dist:eq:cexpand} is justified termwise
for $t$ bounded, but in the limit it is asserted rather than proved.

\subsection{Inter-rupture intervals and the second rupture time}
\label{dist:app:chained-transforms}

Given $k$ founders, the interval $T(k)$ until rupture decomposes over the
holding times: if $G=g$ births precede the catastrophe, the load runs through
$k,k+1,\ldots,k+g$, so
\begin{equation}
T(k)=\sum_{j=k}^{k+G}\xi_j,
\qquad
\xi_j\sim\mathrm{Exp}\lrb{(\lambda+\delta)j}\ \text{independent}.
\label{dist:eq:tkdecomp}
\end{equation}
Averaging over $G\sim\mathrm{Geom}_0(s)$, the Laplace transform is
\begin{equation}
\Lap{T}(k,u)=\E{\ee^{-uT(k)}}
=\sum_{g\ge0}s\rho^g\prod_{j=k}^{k+g}\frac{(\lambda+\delta)j}
{(\lambda+\delta)j+u}.
\label{dist:eq:ltk}
\end{equation}
Writing $u'=u/(\lambda+\delta)$ and unrolling the product with Gamma
functions, the series sums to a hypergeometric form --- a theme that recurs in
\ChPop, whose effective parameters are likewise expressed through
hypergeometric Laplace transforms:
\begin{equation}
\Lap{T}(k,u)=\frac{s\,k}{k+u'}\,\Fhyp\Bigl(k+1,1;k+1+u';\rho\Bigr).
\label{dist:eq:ltkhyp}
\end{equation}
Differentiating \cref{dist:eq:ltk} at $u=0$ gives the exact mean
\begin{equation}
\E{T(k)}=\sum_{m\ge0}\frac{\rho^m}{(\lambda+\delta)(k+m)},
\label{dist:eq:meant}
\end{equation}
which is strictly decreasing in $k$: the more a cell inherits, the faster it
bursts. For $(\lambda,\delta)=(1,1)$ the first four values are
$\E{T(1..4)}\approx0.693,\,0.386,\,0.273,\,0.212$ at fixed founding loads;
the chain intervals $T_k$, whose founding loads are themselves random,
decrease faster still, $0.691,\,0.500,\,0.377,\,0.292$ in simulation.

The chain's dependence structure is visible in the second rupture time. It is
$t_2=t_1+T_2$, but $t_1$ and $T_2$ are \emph{not} independent: both are
functions of $G_1$, the number of births in the first cell. Conditional on
$G_1=g$, however, the two pieces involve disjoint sets of exponential holding
times, so their Laplace transforms factor:
\begin{equation}
\E{\ee^{-ut_2}\mid G_1=g}
=\Pi(1+g,u)\,\Lap{T}(1+g,u),
\qquad
\Pi(m,u)=\prod_{j=1}^{m}\frac{(\lambda+\delta)j}{(\lambda+\delta)j+u}.
\label{dist:eq:t2cond}
\end{equation}
Averaging over $G_1\sim\mathrm{Geom}_0(s)$,
\begin{equation}
\Lap{t_2}(u)
=\sum_{g\ge0} s\rho^g\,\Pi(1+g,u)\,\Lap{T}(1+g,u).
\label{dist:eq:lt2}
\end{equation}
The density itself is computable exactly as a mixture of hypoexponential
densities, since $t_1\mid G_1$ and $T_2\mid(G_1,G_2)$ are sums of independent
exponentials, and it is verified against simulation in
\cref{dist:app:verification}. Moments assemble from the pieces, for instance
$\E{t_2}=\E{T(1)}+\sum_{n\ge1}\pr{r_1=n}\E{T(n)}$.
